\documentclass[border=5pt]{standalone}
\usepackage[T1]{fontenc}
\usepackage{tikz}
\usepackage{luacode}

\definecolor{horseshoe}{HTML}{65780B}
\definecolor{atlantic}{HTML}{466A9F}
\definecolor{garnet}{HTML}{73000A}
\definecolor{rose}{HTML}{CC2E40}
\definecolor{warmgrey}{HTML}{676156}
\definecolor{bk70}{HTML}{5C5C5C}
\definecolor{bk30}{HTML}{C7C7C7}

% Camera icon macro
\newcommand{\drawcameraicon}[1][1]{%
  \begin{scope}[scale=#1]
    \fill[bk70] (-0.18,-0.55) rectangle (0.18,0.55);
    \draw[black, thick] (-0.18,-0.55) rectangle (0.18,0.55);
    \fill[bk30] (0.18,-0.25) rectangle (0.38,0.25);
    \draw[black, thick] (0.18,-0.25) rectangle (0.38,0.25);
    \fill[bk30] (0.38,-0.33) rectangle (0.5,0.33);
    \draw[black, thick] (0.38,-0.33) rectangle (0.5,0.33);
    \fill[rose] (0.0,-0.38) circle (0.07);
    \draw[black, thin] (0.0,-0.38) circle (0.07);
  \end{scope}%
}

% ── CONFIG ──
% camorbitangle set inside tikzpicture for tikzgif
\def\nstops{10}         % number of photo stops around the orbit
\def\snapradius{4}     % how close (in degrees) to a stop to activate rays
\def\objx{2}
\def\objy{0}
\def\objrad{0.6}
\def\orbitradius{2.8}

\begin{document}
\begin{tikzpicture}
  \useasboundingbox (-1.1,-3.1) rectangle (5.1,3.1);
  \draw[horseshoe, thick] (-1.1,-3.1) rectangle (5.1,3.1);

  % Orbit angle from tikzgif \PARAM
  \pgfmathsetmacro{\camorbitangle}{\PARAM}

  % Object
  \fill[warmgrey!25] (\objx,\objy) circle (\objrad);
  \draw[warmgrey, thick] (\objx,\objy) circle (\objrad);

  % Draw stop markers (small dots where camera takes photos)
  \foreach \s in {0,...,9} {
    \pgfmathsetmacro{\stopang}{\s * 360 / \nstops}
    \pgfmathsetmacro{\sx}{\objx + \orbitradius*cos(\stopang)}
    \pgfmathsetmacro{\sy}{\objy + \orbitradius*sin(\stopang)}
    \fill[bk30] (\sx,\sy) circle (0.05);
  }

  % Compute camera position from orbit angle
  \pgfmathsetmacro{\cx}{\objx + \orbitradius*cos(\camorbitangle)}
  \pgfmathsetmacro{\cy}{\objy + \orbitradius*sin(\camorbitangle)}
  \pgfmathsetmacro{\aimang}{atan2(\objy-\cy, \objx-\cx)}

  % Check if camera is near a stop
  \pgfmathsetmacro{\atstop}{0}
  \foreach \s in {0,...,9} {
    \pgfmathsetmacro{\stopang}{\s * 360 / \nstops}
    \pgfmathsetmacro{\diff}{abs(\camorbitangle - \stopang)}
    % Handle wrap-around
    \pgfmathsetmacro{\diff}{min(\diff, 360 - \diff)}
    \ifdim\diff pt<\snapradius pt
      \global\def\atstop{1}
    \fi
  }

  % Lua ray casting (only if at a stop)
\ifnum\atstop=1
\begin{luacode}
local camx = \cx
local camy = \cy
local aim = \aimang * math.pi / 180
local fov = 30 * math.pi / 180
local nrays = 80
local objx, objy, objr = \objx, \objy, \objrad
local max_ray = 4

for i = 0, nrays - 1 do
  local frac = i / (nrays - 1)
  local ang = aim - fov/2 + frac * fov
  local rdx = math.cos(ang)
  local rdy = math.sin(ang)

  local ox = camx - objx
  local oy = camy - objy
  local a = rdx*rdx + rdy*rdy
  local b = 2*(ox*rdx + oy*rdy)
  local c = ox*ox + oy*oy - objr*objr
  local disc = b*b - 4*a*c

  local endx, endy
  local hit = false

  if disc >= 0 then
    local t1 = (-b - math.sqrt(disc)) / (2*a)
    local t2 = (-b + math.sqrt(disc)) / (2*a)
    local t = t1
    if t < 0.01 then t = t2 end
    if t > 0.01 and t < max_ray then
      endx = camx + t * rdx
      endy = camy + t * rdy
      hit = true
    end
  end

  if not hit then
    endx = camx + max_ray * rdx
    endy = camy + max_ray * rdy
  end

  tex.sprint(string.format(
    "\\draw[atlantic, thin, opacity=0.1] (%.4f,%.4f) -- (%.4f,%.4f);",
    camx, camy, endx, endy
  ))
  if hit then
    tex.sprint(string.format(
      "\\fill[atlantic] (%.4f,%.4f) circle (0.03);",
      endx, endy
    ))
  end
end
\end{luacode}
\fi

  % Redraw object on top
  \fill[warmgrey!25] (\objx,\objy) circle (\objrad);
  \draw[warmgrey, thick] (\objx,\objy) circle (\objrad);

  % Camera on top
  \begin{scope}[shift={(\cx,\cy)}, rotate=\aimang]
    \drawcameraicon[0.35]
  \end{scope}

\end{tikzpicture}
\end{document}
